%% SYNTHETIC TEST UNIT generated by writing/tools/make_outline_book.py from FINAL_OUTLINE_v2 -- not book text
\chapter{Collins 减除法的精确形式}\label{app:appN}
本单元为测试用的合成文本（由 FINAL\_OUTLINE\_v2 自动生成），陈述内容从略。

\section{第 1 节（测试）}\label{sec:appN-setting}

\begin{definition}[被积函数、幂次参数、围道与测度]\label{def:appN-integrand}
（测试用合成陈述：内容从略。）
\uses{def:method-power-suppressed}
\end{definition}

\begin{lemma}[幂次压低类的基本性质]\label{lem:appN-suppressed-sets}
（测试用合成陈述：内容从略。）
\uses{def:appN-integrand, def:method-power-suppressed}
\end{lemma}

\begin{remark}[区域集合的两条约定]\label{rem:appN-hard-soft-lines}
（测试用合成陈述：内容从略。）
% uses: none
\end{remark}

\begin{lemma}[邻域指数的可行选择]\label{lem:appN-beta}
（测试用合成陈述：内容从略。）
\uses{def:method-neighbourhoods}
\end{lemma}

\begin{example}[两圈平面顶点的九个领头区域]\label{ex:appN-vertex-regions}
（测试用合成陈述：内容从略。）
\uses{def:method-regions, lem:method-meets, lem:appN-beta}
\end{example}

\begin{assumption}[(GA1) 围道 \tagA]\label{asm:appN-contour}
（测试用合成陈述：内容从略。）
地位：测试状态句。
文献：测试。
被使用于：\cref{cor:method-removing-eps}、\cref{thm:method-subtraction}、\cref{prop:disall-two-routes}、\cref{thm:disall-renorm-compat}、\cref{thm:appN-subtraction}、\cref{prop:appN-sia}、\cref{thm:appN-structure}。
解除途径：测试解除途径。
\uses{def:appN-integrand, def:method-subtraction}
\end{assumption}
\begin{proofsketch}测试。\end{proofsketch}

\begin{definition}[(GA2) 对逐线交封闭的区域集合]\label{def:appN-meet-closed}
（测试用合成陈述：内容从略。）
\uses{def:method-regions, lem:method-meets}
\end{definition}

\section{第 2 节（测试）}\label{sec:appN-inputs}

\begin{assumption}[(GA3) 幂次计数包络 \tagA]\label{asm:appN-envelope}
（测试用合成陈述：内容从略。）
地位：测试状态句。
文献：测试。
被使用于：\cref{lem:method-local-accuracy}、\cref{thm:method-template}、\cref{prop:disall-two-routes}、\cref{thm:disall-factorization}、\cref{rem:appE-actual-bound}、\cref{lem:appN-error-table}、\cref{prop:appN-sia}、\cref{thm:appN-structure}、\cref{prop:appO-one-sided}、\cref{thm:appO-main}、\cref{cor:appO-interface}、\cref{prop:appO-tmd}。
解除途径：测试解除途径。
\uses{def:method-neighbourhoods, def:method-collinear-approx, def:method-soft-approx}
\end{assumption}
\begin{proofsketch}测试。\end{proofsketch}

\begin{assumption}[(GA4) 分离性 \tagA]\label{asm:appN-separation}
（测试用合成陈述：内容从略。）
地位：测试状态句。
文献：测试。
被使用于：\cref{lem:method-local-accuracy}、\cref{thm:method-template}、\cref{prop:disall-two-routes}、\cref{thm:disall-factorization}、\cref{lem:appN-error-table}、\cref{prop:appN-sia}、\cref{prop:appN-tmd-window}、\cref{thm:appN-structure}、\cref{prop:appO-one-sided}、\cref{thm:appO-main}、\cref{cor:appO-interface}、\cref{prop:appO-tmd}。
解除途径：测试解除途径。
\uses{def:method-neighbourhoods, def:ir-glauber}
\end{assumption}
\begin{proofsketch}测试。\end{proofsketch}

\begin{remark}[各输入的地位]\label{rem:appN-inputs-status}
（测试用合成陈述：内容从略。）
% uses: none
\end{remark}

\begin{theorem}[减除定理（条件形式） \tagR]\label{thm:appN-subtraction}
（$R_{\mathrm{cond}}$：在所列编号假设之下证明完整；结论的地位不高于其中最弱的假设。）
前提：\cref{asm:appN-contour}。
（测试用合成陈述：内容从略。）
（本命题只涉及部分子态。）
% errtype: E-N
\uses{asm:appN-contour, def:method-collins-conditions, def:method-subtraction, lem:method-covering, lem:appN-suppressed-sets}
\end{theorem}

\begin{definition}[纤维化密度、可检验条件 (FR)(FI) 与误差类型 E-N$^\varphi$]\label{def:appN-fibered}
（测试用合成陈述：内容从略。）
% errtype: E-Nphi
\uses{def:appN-integrand, def:method-subtraction, def:method-collins-conditions}
\end{definition}

\begin{theorem}[纤维化减除定理]\label{thm:appN-fibered-subtraction}
（测试用合成陈述：内容从略。）
\uses{def:appN-fibered, def:method-subtraction, def:method-collins-conditions}
\end{theorem}

\begin{lemma}[纯分母乘积的纤维化–协变等同]\label{lem:appN-identification}
（测试用合成陈述：内容从略。）
\uses{def:appN-fibered}
\end{lemma}

\section{第 3 节（测试）}\label{sec:appN-theorem}

\begin{lemma}[热带/线性规划证书]\label{lem:appN-lp-certificate}
（测试用合成陈述：内容从略。）
\uses{def:appN-fibered}
\end{lemma}

\begin{lemma}[光锥能量分母的定号性（标量三次模型、无质量内线） \tagR]\label{lem:appN-lc-positivity}
（$R_{\mathrm{PT}}$：在重整化微扰论中对红外正规化的部分子态严格。）
（测试用合成陈述：内容从略。）
（本命题只涉及部分子态。）
\uses{def:appN-fibered}
\end{lemma}

\begin{assumption}[光锥序表示与协变值的等同（LCE） \tagA]\label{asm:appN-lce}
（测试用合成陈述：内容从略。）
地位：测试状态句。
文献：测试。
被使用于：\cref{cor:appN-cs3-fibered}。
解除途径：测试解除途径。
\uses{def:appN-fibered, lem:appN-lc-positivity}
\end{assumption}
\begin{proofsketch}测试。\end{proofsketch}

\begin{corollary}[红外有限图的 (CS3$^\varphi$) \tagR]\label{cor:appN-cs3-fibered}
（$R_{\mathrm{cond}}$：在所列编号假设之下证明完整；结论的地位不高于其中最弱的假设。）
前提：\cref{asm:appN-lce}。
（测试用合成陈述：内容从略。）
（本命题只涉及部分子态。）
\uses{asm:appN-lce, lem:appN-lc-positivity}
\end{corollary}

\begin{remark}[与 BPHZ 森林公式的正确类比]\label{rem:appN-forest}
（测试用合成陈述：内容从略。）
% uses: none
\end{remark}

\begin{assumption}[发散链项的取值规则]\label{asm:appN-values}
（测试用合成陈述：内容从略。）
地位：测试状态句。
文献：测试。
被使用于：\cref{thm:method-template}、\cref{prop:appN-thrust-zerobin}、\cref{prop:appN-thrust}、\cref{thm:appN-structure}、\cref{prop:appN-triangle-exact}、\cref{prop:appN-twoloop-remainder}。
解除途径：测试解除途径。
\uses{def:method-subtraction, prop:method-chains}
\end{assumption}

\begin{assumption}[一般图的 Collins 条件 \tagA]\label{asm:appN-collins}
（测试用合成陈述：内容从略。）
地位：测试状态句。
文献：测试。
被使用于：\cref{thm:method-subtraction}、\cref{thm:method-template}、\cref{prop:disall-two-routes}、\cref{prop:disall-h5}、\cref{thm:disall-factorization}、\cref{thm:dy-factorization}、\cref{thm:tmd-factorization}、\cref{prop:excl-soft-cancel}、\cref{thm:excl-meson-L}、\cref{prop:appN-sia}、\cref{prop:appN-thrust}、\cref{prop:appN-dy}、\cref{prop:appN-tmd-window}、\cref{thm:appN-structure}。
解除途径：测试解除途径。
\uses{def:method-collins-conditions, def:appN-meet-closed}
\end{assumption}
\begin{proofsketch}测试。\end{proofsketch}

\section{第 4 节（测试）}\label{sec:appN-approximators}

\begin{definition}[TMD 近似子与快度正规化]\label{def:appN-tmd-approx}
（测试用合成陈述：内容从略。）
\uses{def:method-collinear-approx, def:method-soft-approx}
\end{definition}

\begin{definition}[遍举近似子与分布振幅投影]\label{def:appN-excl-approx}
（测试用合成陈述：内容从略。）
\uses{def:method-collinear-approx, def:method-soft-approx}
\end{definition}

\begin{lemma}[近似子的代数]\label{lem:appN-algebra}
（测试用合成陈述：内容从略。）
\uses{def:method-collinear-approx, def:method-soft-approx, lem:method-meets}
\end{lemma}

\begin{remark}[(Abs) 的一个未被覆盖的情形]\label{rem:appN-abs-edge}
（测试用合成陈述：内容从略。）
% uses: none
\end{remark}

\begin{remark}[没有适配动量分配时]\label{rem:appN-routing}
（测试用合成陈述：内容从略。）
% uses: none
\end{remark}

\begin{lemma}[误差表各行的证明 \tagR]\label{lem:appN-error-table}
（$R_{\mathrm{cond}}$：在所列编号假设之下证明完整；结论的地位不高于其中最弱的假设。）
前提：\cref{asm:appN-separation}、\cref{asm:appN-envelope}。
（测试用合成陈述：内容从略。）
（本命题只涉及部分子态。）
% errtype: E-P
\uses{asm:appN-separation, asm:appN-envelope, lem:appN-beta, def:method-collinear-approx, def:method-soft-approx, def:appN-tmd-approx, def:appN-excl-approx}
\end{lemma}

\begin{assumption}[Ward 恒等式的乘积结构与减除的重组 \tagA]\label{asm:appN-ward-product}
（测试用合成陈述：内容从略。）
地位：测试状态句。
文献：测试。
被使用于：\cref{thm:method-template}、\cref{prop:disall-factorized-subtraction}、\cref{thm:disall-factorization}、\cref{prop:tmd-zero-bin-soft}、\cref{prop:excl-soft-cancel}、\cref{prop:appN-sia}、\cref{prop:appN-thrust}、\cref{prop:appN-dy}、\cref{thm:appN-structure}、\cref{prop:appO-past-wilson}、\cref{thm:appO-main}、\cref{cor:appO-interface}、\cref{prop:appO-tmd}。
解除途径：测试解除途径。
\uses{lem:method-ward-multi, prop:method-collinear-wilson, prop:method-soft-wilson, def:method-subtraction}
\end{assumption}
\begin{proofsketch}测试。\end{proofsketch}

\section{第 5 节（测试）}\label{sec:appN-processes}

\begin{assumption}[多定域光线算符的 BRST 分类 \tagC]\label{asm:appN-brst-nonlocal}
（测试用合成陈述：内容从略。）
地位：测试状态句。
文献：测试。
被使用于：\cref{thm:appN-structure}。
解除途径：测试解除途径。
\uses{def:appM-alien, prop:appJ-local-operators}
\end{assumption}

\begin{assumption}[割图求和的公共表示 \tagA]\label{asm:appN-cutsum}
（测试用合成陈述：内容从略。）
地位：测试状态句。
文献：测试。
被使用于：\cref{prop:method-fsi}、\cref{thm:method-template}、\cref{prop:disall-two-routes}、\cref{thm:disall-factorization}、\cref{thm:dy-factorization}、\cref{prop:appN-sia}、\cref{prop:appN-thrust}、\cref{prop:appN-dy}、\cref{prop:appN-tmd-window}、\cref{thm:appN-structure}。
解除途径：测试解除途径。
\uses{def:appN-integrand, def:prelim-cut}
\end{assumption}
\begin{proofsketch}测试。\end{proofsketch}

\begin{assumption}[被积函数层次的紫外重整化 \tagA]\label{asm:appN-uv}
（测试用合成陈述：内容从略。）
地位：测试状态句。
文献：测试。
被使用于：\cref{thm:method-subtraction}、\cref{thm:method-template}、\cref{prop:disall-factorized-subtraction}、\cref{thm:disall-factorization}、\cref{thm:appN-structure}。
解除途径：测试解除途径。
\uses{def:appN-integrand, thm:prelim-composite-renorm}
\end{assumption}
\begin{proofsketch}测试。\end{proofsketch}

\begin{assumption}[共线因子的乘法重整化（维数正规化路线） \tagA]\label{asm:appN-multren}
（测试用合成陈述：内容从略。）
地位：测试状态句。
文献：测试。
被使用于：\cref{prop:pdf-timelike-dglap}、\cref{thm:method-template}、\cref{thm:disall-renorm-compat}、\cref{prop:disall-h5}、\cref{thm:disall-factorization}、\cref{thm:coll-sia}、\cref{asm:appN-ir-regulator}、\cref{thm:appN-structure}。
解除途径：测试解除途径。
\uses{thm:pdf-moments-local, prop:pdf-partonic-dr}
\end{assumption}
\begin{proofsketch}测试。\end{proofsketch}

\begin{assumption}[保持 Ward 恒等式的部分子靶红外正规化 \tagC]\label{asm:appN-ir-regulator}
（测试用合成陈述：内容从略。）
地位：测试状态句。
文献：测试。
被使用于：\cref{thm:disall-renorm-compat}。
解除途径：测试解除途径。
\uses{asm:appN-multren}
\end{assumption}

\begin{proposition}[$e^+e^-\to hX$ 的减除结构 \tagR]\label{prop:appN-sia}
（$R_{\mathrm{cond}}$：在所列编号假设之下证明完整；结论的地位不高于其中最弱的假设。）
前提：\cref{asm:appN-collins}、\cref{asm:appN-cutsum}、\cref{asm:appN-contour}、\cref{asm:appN-envelope}、\cref{asm:appN-separation}、\cref{asm:appN-ward-product}、\cref{lem:coll-unobserved-jet}。
（测试用合成陈述：内容从略。）
（本命题只涉及部分子态。）
% errtype: E-N
\uses{asm:appN-collins, asm:appN-cutsum, asm:appN-contour, asm:appN-envelope, asm:appN-separation, asm:appN-ward-product, thm:appN-subtraction, lem:coll-final-poles, lem:coll-unobserved-jet, lem:method-soft-unitarity, lem:method-i0-compat}
\end{proposition}

\begin{proposition}[推力的零箱积分无标度 \tagR]\label{prop:appN-thrust-zerobin}
（$R_{\mathrm{cond}}$：在所列编号假设之下证明完整；结论的地位不高于其中最弱的假设。）
前提：\cref{asm:appN-values}。
（测试用合成陈述：内容从略。）
（本命题只涉及部分子态。）
\uses{asm:appN-values, def:method-collinear-approx, def:method-soft-approx}
\end{proposition}

\section{第 6 节（测试）}\label{sec:appN-exclusive}

\begin{proposition}[双喷注事件形状（推力）的减除结构 \tagR]\label{prop:appN-thrust}
（$R_{\mathrm{cond}}$：在所列编号假设之下证明完整；结论的地位不高于其中最弱的假设。）
前提：\cref{asm:appN-collins}、\cref{asm:appN-cutsum}、\cref{asm:appN-ward-product}、\cref{asm:appN-values}。
（测试用合成陈述：内容从略。）
（本命题只涉及部分子态。）
\uses{asm:appN-collins, asm:appN-cutsum, asm:appN-ward-product, asm:appN-values, thm:appN-subtraction, prop:appN-thrust-zerobin, lem:appN-error-table}
\end{proposition}

\begin{proposition}[DY 共线因子化的条件结构 \tagR]\label{prop:appN-dy}
（$R_{\mathrm{cond}}$：在所列编号假设之下证明完整；结论的地位不高于其中最弱的假设。）
前提：\cref{asm:appN-collins}、\cref{asm:appN-cutsum}、\cref{asm:appN-ward-product}、\cref{asm:appO-rewritten-cs}、\cref{asm:appO-deformation}、\cref{cor:appO-interface}、\cref{prop:method-collinear-wilson}。
（测试用合成陈述：内容从略。）
（本命题只涉及部分子态。）
% errtype: E-W
\uses{asm:appN-collins, asm:appN-cutsum, asm:appN-ward-product, asm:appO-rewritten-cs, asm:appO-deformation, cor:appO-interface, lem:method-soft-unitarity, thm:appN-subtraction, prop:method-collinear-wilson, lem:method-i0-compat}
\end{proposition}

\begin{assumption}[TMD 输入 \tagA]\label{asm:appN-tmd}
（测试用合成陈述：内容从略。）
地位：测试状态句。
文献：测试。
被使用于：\cref{prop:tmd-zero-bin-soft}、\cref{thm:tmd-factorization}、\cref{prop:tmd-evolution}、\cref{prop:appN-tmd-window}、\cref{thm:appN-structure}。
解除途径：测试解除途径。
\uses{def:appN-tmd-approx, def:tmd-unsub, def:tmd-soft}
\end{assumption}
\begin{proofsketch}测试。\end{proofsketch}

\begin{remark}[类空正规子线的自能]\label{rem:appN-spacelike-selfenergy}
（测试用合成陈述：内容从略。）
% uses: none
\end{remark}

\begin{lemma}[一阶零箱吸收关系 \tagR]\label{lem:appN-absorption}
（$R_{\mathrm{PT}}$：在重整化微扰论中对红外正规化的部分子态严格。）
（测试用合成陈述：内容从略。）
（本命题只涉及部分子态。）
\uses{def:appN-tmd-approx, def:method-subtraction}
\end{lemma}

\begin{proposition}[TMD 快度窗口中的减除定理 \tagR]\label{prop:appN-tmd-window}
（$R_{\mathrm{cond}}$：在所列编号假设之下证明完整；结论的地位不高于其中最弱的假设。）
前提：\cref{asm:appN-collins}、\cref{asm:appN-cutsum}、\cref{asm:appN-tmd}、\cref{asm:appN-separation}、\cref{prop:appO-tmd}。
（测试用合成陈述：内容从略。）
（本命题只涉及部分子态。）
\uses{asm:appN-collins, asm:appN-cutsum, asm:appN-tmd, asm:appN-separation, prop:appO-tmd, thm:appN-subtraction, lem:appN-error-table}
\end{proposition}

\begin{assumption}[端点条件 \tagA]\label{asm:appN-endpoint}
（测试用合成陈述：内容从略。）
地位：测试状态句。
文献：测试。
被使用于：\cref{prop:excl-soft-cancel}、\cref{thm:excl-pion-ff}、\cref{thm:excl-meson-L}、\cref{thm:appN-structure}。
解除途径：测试解除途径。
\uses{def:appN-excl-approx, def:method-collins-conditions, def:excl-da}
\end{assumption}
\begin{proofsketch}测试。\end{proofsketch}

\section{第 7 节（测试）}\label{sec:appN-examples}

\begin{theorem}[由减除法得到的因子化结构（条件形式） \tagR]\label{thm:appN-structure}
（$R_{\mathrm{cond}}$：在所列编号假设之下证明完整；结论的地位不高于其中最弱的假设。）
前提：\cref{asm:appN-contour}、\cref{asm:appN-envelope}、\cref{asm:appN-separation}、\cref{asm:appN-collins}、\cref{asm:appN-uv}、\cref{asm:appN-multren}、\cref{asm:appN-ward-product}、\cref{asm:appN-brst-nonlocal}、\cref{asm:appN-cutsum}、\cref{asm:appN-values}、\cref{cor:appO-interface}、\cref{asm:appN-tmd}、\cref{prop:appO-tmd}、\cref{asm:appN-endpoint}。
（测试用合成陈述：内容从略。）
（本命题只涉及部分子态。）
\uses{asm:appN-contour, asm:appN-envelope, asm:appN-separation, asm:appN-collins, asm:appN-uv, asm:appN-multren, asm:appN-ward-product, asm:appN-brst-nonlocal, asm:appN-cutsum, asm:appN-values, asm:appN-tmd, asm:appN-endpoint, cor:appO-interface, prop:appO-tmd, thm:appN-subtraction, prop:disall-two-routes, prop:disall-factorized-subtraction, prop:appN-sia, prop:appN-thrust, prop:appN-dy, prop:appN-tmd-window, prop:tmd-zero-bin-soft, lem:tmd-sqrt-identity, prop:excl-soft-cancel, lem:method-soft-unitarity}
\end{theorem}

\begin{proposition}[单圈离壳三角图：Collins 减除的完整计算与精确余项 \tagR]\label{prop:appN-triangle-exact}
（$R_{\mathrm{PT}}$：在重整化微扰论中对红外正规化的部分子态严格。）
（测试用合成陈述：内容从略。）
（本命题只涉及部分子态。）
\uses{prop:method-product, lem:appN-algebra, asm:appN-values}
\end{proposition}

\begin{proposition}[两级梯形图：嵌套区域与余项的乘积结构 \tagR]\label{prop:appN-ladder}
（$R_{\mathrm{PT}}$：在重整化微扰论中对红外正规化的部分子态严格。）
（测试用合成陈述：内容从略。）
（本命题只涉及部分子态。）
\uses{prop:method-product, lem:appN-algebra}
\end{proposition}

\begin{proposition}[单圈离壳三角图：纤维化 Collins 条件的完整验证与积分减除定理 \tagR]\label{prop:appN-triangle-cs}
（$R_{\mathrm{PT}}$：在重整化微扰论中对红外正规化的部分子态严格。）
（测试用合成陈述：内容从略。）
（本命题只涉及部分子态。）
\uses{prop:appN-triangle-exact, def:appN-fibered, thm:appN-fibered-subtraction, lem:appN-lp-certificate}
\end{proposition}

\begin{proposition}[$\varphi^3_6$ 中 DIS 前向振幅：全部单圈图与单梯级梯形图 \tagR]\label{prop:appN-phi36-dis}
（$R_{\mathrm{PT}}$：在重整化微扰论中对红外正规化的部分子态严格。）
（测试用合成陈述：内容从略。）
（本命题只涉及部分子态。）
\uses{prop:appN-ladder, def:appN-fibered, thm:appN-fibered-subtraction}
\end{proposition}

\begin{proposition}[阿贝尔有质量光子：类光近似子的精确余项与快度正规子 \tagR]\label{prop:appN-massive-photon}
（$R_{\mathrm{PT}}$：在重整化微扰论中对红外正规化的部分子态严格。）
（测试用合成陈述：内容从略。）
（本命题只涉及部分子态。）
\uses{def:appN-fibered, def:appN-tmd-approx}
\end{proposition}

\begin{proposition}[两圈平面顶点：区域代数与字面约定下 (Abs) 的失效 \tagR]\label{prop:appN-twoloop-algebra}
（$R_{\mathrm{PT}}$：在重整化微扰论中对红外正规化的部分子态严格。）
（测试用合成陈述：内容从略。）
（本命题只涉及部分子态。）
\uses{ex:appN-vertex-regions, lem:appN-algebra}
\end{proposition}

\section{第 8 节（测试）}\label{sec:appN-status}

\begin{assumption}[两圈图样积分的辅助解析正规化 \tagA]\label{asm:appN-analytic-reg}
（测试用合成陈述：内容从略。）
地位：测试状态句。
文献：测试。
被使用于：\cref{prop:appN-twoloop-remainder}。
解除途径：测试解除途径。
\uses{prop:appN-twoloop-algebra}
\end{assumption}
\begin{proofsketch}测试。\end{proofsketch}

\begin{proposition}[两圈平面顶点：Collins 项与区域积分的等同及积分余项 \tagR]\label{prop:appN-twoloop-remainder}
（$R_{\mathrm{cond}}$：在所列编号假设之下证明完整；结论的地位不高于其中最弱的假设。）
前提：\cref{asm:appN-analytic-reg}、\cref{asm:appN-values}。
（测试用合成陈述：内容从略。）
（本命题只涉及部分子态。）
\uses{asm:appN-analytic-reg, asm:appN-values, prop:appN-twoloop-algebra}
\end{proposition}

\begin{remark}[超可重整模型 (M1) 的红外结构：全阶例子需要的模型]\label{rem:appN-m1-infrared}
（测试用合成陈述：内容从略。）
% uses: none
\end{remark}

\begin{remark}[全阶的纤维化减除：条件形式、范围与定位]\label{rem:appN-fibered-allorders}
（测试用合成陈述：内容从略。）
% uses: none
\end{remark}
